% --------------------------------------------------------------------------
\section{Multi-Attack Training as a Domain Problem}
\label{sec:ch4_problem_formulation}

\subsection{Setup and Notation}

Let
\begin{equation}
    \mathcal{D}_{\mathrm{train}}
    =
    \{(\mathbf{x}_i, y_i)\}_{i=1}^{n}
\end{equation}
denote the supervised training dataset, where \(\mathbf{x}_i \in \mathcal{X}\) is an
input image and \(y_i \in \{1,\ldots,C\}\) is its class label. Let \(f_\theta\) be a
classifier parameterized by \(\theta\), and let
\(\ell(f_\theta(\mathbf{x}),y)\) denote the cross-entropy loss.

We denote the logits produced by the model as
\[
    \mathbf{z} = f_\theta(\mathbf{x}) \in \mathbb{R}^{C},
\]
and the penultimate-layer representation as
\[
    \mathbf{h} = \phi_\theta(\mathbf{x}) \in \mathbb{R}^{d_h}.
\]
The representation \(\phi_\theta(\cdot)\) is later used for cluster discovery.

We define a set of source attack domains
\begin{equation}
    \mathcal{A}_{\mathrm{src}}
    =
    \{A_1,A_2,\ldots,A_M\},
\end{equation}
where each attack operator \(A_m\) is characterized by a threat model, perturbation
budget, optimization procedure, number of inner steps, and loss objective. Given a
clean sample \((\mathbf{x}_i,y_i)\), the adversarial example produced by source attack
\(A_m\) is
\begin{equation}
    \mathbf{x}_i^{(m)}
    =
    A_m(f_\theta,\mathbf{x}_i,y_i).
\end{equation}

In the main experiments, we use two source attacks: a PGD-\(\ell_\infty\) attack and
a DDN-\(\ell_2\) attack. This choice exposes the model to two common adversarial norm
families during training. Additional attacks are reserved as heldout evaluation attacks
and are never used for optimization.

\subsection{Multi-Attack Risk and Its Limitations}

The empirical adversarial risk under source attack \(A_m\) is
\begin{equation}
    R_m(\theta)
    =
    \frac{1}{n}
    \sum_{i=1}^{n}
    \ell\!\left(
        f_\theta\!\left(\mathbf{x}_i^{(m)}\right),
        y_i
    \right).
\end{equation}

A standard multi-attack objective minimizes the average risk across source attacks:
\begin{equation}
    \min_\theta
    \frac{1}{M}
    \sum_{m=1}^{M}
    R_m(\theta).
    \label{eq:avg_risk}
\end{equation}

This objective improves attack diversity compared with single-attack adversarial
training. However, it treats all source attacks and all adversarial examples with equal
importance. This creates two limitations.

First, if one source attack consistently induces higher loss than another, equal
weighting may underemphasize the harder attack. Second, adversarial examples within
the same attack are not homogeneous. Two samples generated by the same algorithm can
differ substantially in classification margin, representation, and optimization
direction. These within-attack differences are not captured by
Equation~\eqref{eq:avg_risk}.

The central goal of this chapter is therefore to construct adaptive objectives that can
identify and upweight hard groups while preserving the stability of the uniform
multi-attack training signal.
